\chapter[Universal Tensor in Sobolev Norm]{The Universal Tensor in the Sobolev Norm}
\label{chap:universal-tensor-sobolev}

The construction implies a tensor before it implies a continuum. That order is important. A tensor is the object
that knows how the components of the record pair, contract, and transform. A continuum field is a later reading,
a smooth shadow drawn through refinements. The finite apparatus therefore begins with a universal tensor row
system and only then asks which Sobolev norm is appropriate for the experimental floor.

\section{The Two-Term Residual}

The capstone language of the construction separates two terms. The zeroth term is the tensor value: the carrier
of geometry, boundary, and frame. The first term is the variation: the derivative face that, in the quantum
regime, is read through the Dirac operator. In the weak experiment these two terms meet as a residual:
\[
  R_U(\psi) =
  \mathcal D_A\psi + \mathcal T_U\psi - f .
\]
The subscript \(U\) records that this is not an arbitrary tensor. It is the universal tensor as the apparatus
can express it at the chosen finite floor. The Dirac term supplies first-order spinor transport. The tensor term
supplies the background contraction. The load supplies the experiment. A solution is not a state that makes the
symbols pretty. It is a state whose residual cannot be detected by any admissible test at the chosen Sobolev
resolution.

The Lean file represents this two-term residual exactly:

\begin{quote}\small
\begin{verbatim}
def matrixApply (rows : List Vec) (state : Vec) : Vec :=
  rows.map (fun row => rowApply row state)

def rawTensorResidual
    (U : UniversalTensor) (state : Vec) : Vec :=
  sub
    (add (matrixApply U.diracRows state)
      (matrixApply U.tensorRows state))
    U.load
\end{verbatim}
\end{quote}

The important word is \(\texttt{rows}\). A Galerkin method begins by choosing a finite test and trial space.
Once that choice is made, the operator is a table of exact row functionals. The Dirac operator has rows. The
tensor has rows. The load has rows. The weak problem becomes a finite residual ledger.

\section{The Sobolev Reader}

The Sobolev norm says how the residual is read. In analysis it may be
\[
  \|r\|_{H^s}^{2}
  =
  \sum_{|\alpha|\le s}\|D^\alpha r\|_{L^2}^{2},
\]
or an equivalent form built from an elliptic operator. The details depend on the domain, boundary conditions,
and regularity. The apparatus does not need to pretend that every experiment uses the same \(s\). It needs to
record the Sobolev reader explicitly.

In the finite companion, that reader is another row system:

\begin{quote}\small
\begin{verbatim}
def sobolevResidual
    (U : UniversalTensor) (state : Vec) : Vec :=
  matrixApply U.sobolevRows (rawTensorResidual U state)

def sobolevNormSq (U : UniversalTensor) (state : Vec) : Int :=
  let r := sobolevResidual U state
  dot r r
\end{verbatim}
\end{quote}

The row system may be the identity, a stiffness matrix, a mass-plus-stiffness matrix, or a boundary-weighted
reader. The point is not which Sobolev norm is chosen in the toy file. The point is that the Sobolev norm is not
an afterthought. It is part of the experimental apparatus. Changing it changes what counts as small, what counts
as zero, and what the weak equation asks the state to satisfy.

\section{Principle of Least Activity}

The principle is least activity, not least decoration. The activity is
\[
  \mathcal A_U(\psi) =
  \frac12\|R_U(\psi)\|_{H^s}^{2}.
\]
For a linearized residual, the first Fréchet variation against a test state \(\eta\) is the Sobolev pairing of
the residual with that test:
\[
  D\mathcal A_U(\psi)[\eta]
  =
  \langle R_U(\psi), \eta\rangle_{H^s}.
\]
The weak equation is
\[
  D\mathcal A_U(\psi)[\eta] = 0
  \qquad
  \hbox{for every admissible test }\eta .
\]
This is the sentence the user gave, with the roles separated: the weak solution to the principle of least
activity is the state whose first Fréchet variation of the universal-tensor residual vanishes in the Sobolev
norm. The solution is not the derivative. The solution is what makes the derivative vanish.

The finite file writes the pairing as:

\begin{quote}\small
\begin{verbatim}
def firstFrechetVariation
    (U : UniversalTensor) (state test : Vec) : Int :=
  dot (sobolevResidual U state) test

def WeakLeastActivity
    (U : UniversalTensor) (state : Vec) : Prop :=
  forall test, firstFrechetVariation U state test = 0
\end{verbatim}
\end{quote}

All of the analysis has been squeezed to its exact computational core. A residual is built. A Sobolev reader
acts on it. A test vector pairs with it. The weak solution is the state for which every such pairing vanishes.
